\documentclass[book.tex]{subfiles}

\begin{document}
\chapter{Quantum Field Theory Basics}

\label{chap:qft_basics}
\noindent\rule{11cm}{0.4pt} \\
\textbf{Prerequisites:} \autoref{chap:schrodinger}, \autoref{chap:dirac_equation} \\
\textbf{Difficulty Level:} ** \\
\noindent\rule{11cm}{0.4pt}
\vspace{5mm}

Quantum field theory builds machinery to handle relativistic quantum mechanical systems with an unbounded number of particles. 
This chapter will introduce considerable mathematical machinery, but we will aim to motivate this machinery by tying back to examples you've already studied at depth through this book.


\section{Relativistic Quantum Mechanics}

The standard form of Schrodinger's equation we have seen thus far is not relativistically invariant. Quantum field theory arose from attempts to modify Schrodinger's equation to become relativistically invariant by modifying the base form of Schrodinger's equation to respect Lorentz invariance. The Klein-Gordon equation provides a simple first relativistically invariant variant of Schrodinger's equation for spinless particles without interactions.

The Dirac equation extends the base Klein-Gordon equation to add spin, necessary to model particles such as electrons, but requires the introduction of more sophisticated machinery in terms of spinors. We will return to this equation shortly.

\section{Quantum Field Theory}

The basic idea of second quantization is that fields themselves need to be promoted to operator quantities in order to model quantum field-theoretic systems. This is called second quantization as opposed to first quantization, which is the procedure used for classical quantities like position and momentum. 


%\textcolor{red}{TODO: Define here or in mathematical preliminaries chapter}

\section{Quantum Field Theory as a Physical Theory}
%\textcolor{red}{\url{https://statweb.stanford.edu/~souravc/qft-lectures-combined.pdf} <- Incorporate more from this reference}


States $\mathcal{S}$ are given by Hilbert space $\mathcal{H}$ of wave functions as we saw earlier for quantum mechanics, but we will need to use a larger Hilbert space. In particular, we will choose the Fock space. Let's quickly motivate this construction.  The suitable Hilbert space for particle is $\mathcal{H} = L^2(\mathbb{R}^3)$. The Hilbert space for $n$ particles is the tensor product $\mathcal{H}^{\otimes n}$. Let $S_+$ be the tensor symmetrization operator. Then the bosonic Fock space is
    
\begin{align} F_+(\mathcal{H}) = \overline{\bigoplus_{n=0}^\infty S_+(\mathcal{H}^{\otimes n}}) \end{align}
    
The bar above indicates the Hilbert space completion. This definition is a little complicated to digest, but you can think of it as allowing for quantum states of arbitrarily many particles.
    
What then is a field? In mathematical language, a field $\phi$ is an operator-valued distribution defined on $F_+(\mathcal{H})$. That is, $\phi$ maps functions in $F_+(\mathcal{H})$ to operators. If we squint we can pretend that points in $\mathbb{R}^3$ are points in $F_+(\mathcal{H})$. (Pretend the Dirac delta $\delta_x$ at $x \in \mathbb{R}^3$ is in $\mathcal{H} = L^2(\mathbb{R}^n)$. Then the field $\phi$ maps $\delta_x$ to an operator. That is, there is an operator for every point of space!)

For symmetries $\mathcal{G}$, unlike in nonrelativistic quantum mechanics, it is more convenient to define a quantum field theory by its Lagrangian. Here we describe a simple field theory with quartic self-interaction
    
\begin{align} \mathcal{L} = \frac{1}{2} (\partial_{\mu} \phi) (\partial^{\mu} \phi) - \frac{1}{2} m^2 \phi^2 - \frac{\lambda}{4!} \phi^4 \end{align}
    
Here $\mathcal{L}$ is another operator-valued distribution. As a mathematical note, the Hilbert space $F_+(\mathcal{H})$ we defined above isn't right. You can pretend we still live in this space, but it's an open mathematical problem to find the "correct" space.

We define the time evolution operator $U(t_2, t_1)$ to be the operator such that
\begin{align} | \psi(t_2) \rangle = U(t_2, t_1) |\psi(t_1) \rangle \end{align}
This time evolution operator can be expanded in terms of the Dyson series
\begin{align} 
U(t_2, t_1) &= T\left [ \exp \left ( -i \int_{t_1}^{t_2} dt H_I(t) \right ) \right ]
\end{align}

where $T$ is the time ordering operator and $H_I$ is the interaction Hamiltonian. Parameters are given by

\begin{align}\mathcal{W} = \{m, \lambda, |\psi(0)\rangle \}\end{align}
where $m$, $\lambda$ are the "bare mass" and "bare coupling constant" and $|\psi(0)\rangle$ is the initial condition of the wave function.

The domain of Applicability $\mathcal{C}$ is set by constraints
Mass $m_r < M$ where $m_r$ is the physically measureable mass and $M$ is some ceiling parameter. Hyperparameters $\mathcal{HP}$ are $M$, the mass limit, and noise parameter $\epsilon$


\begin{center}
\begin{tabular}{|c|c|}
    \hline
    States & $\mathcal{S} = F_+(\mathcal{H}) $ \\
    \hline
    Symmetries & $\mathcal{G} = \{\}$\\
    \hline
    Operators & $\mathcal{O} = \{ T \left [ \exp\left ( -i \int_{t_1}^{t_2} dt H_I(t) \right ) \right ]\} $\\
    \hline
    Parameters & $\mathcal{W} = \{m, \lambda, |\psi(0)\rangle \}$ \\
    \hline
    Constraints & $\mathcal{C} =  \{m_r < M\}$ \\
    \hline
    Hyperparameters & $\mathcal{HP} = \{M, \epsilon\}$  \\
    \hline
    Observables & $\mathcal{B} = \{f(x) = x + \mathcal{N}(b, \sigma) \}$  \\
    \hline
\end{tabular}
\label{phs:theory:qft}
\end{center}
\section{Mathematical Preliminaries}

In this section, we introduce a few additional pieces of mathematical infrastructure that we will find useful in subsequent chapters.

\subsection{Scattering Matrix}
Physical observables in quantum field theory are generally written in terms of the scattering matrix $S$
\begin{align} S = \lim_{t_2 \to \infty} \lim_{t_1 \to -\infty} U(t_2, t_1) \end{align}
This is called a "matrix", since the intuition is that term $\langle \psi_i | S | \psi_j \rangle$ captures the probability of "basis state" $\psi_j$ transitioning to $\psi_i$ over infinite time. If the input and output states are well chosen, this term can be estimated from experiment with some noise $\textrm{Normal}(0, \epsilon)$
\subsection{Second Quantization}
Let's define the number operator $N$
\begin{equation} 
N = \sum_i \hat{a}_i^{\dagger} \hat{a}_i
\end{equation}
The number operator counts the number of particles in a system. Here $\hat{a}_i^\dagger$ is the creation operator and $\hat{a}_i$ is the annihilation operator. 

\subsection{Operator Valued Distributions}
Quantum fields are typically mathematically modeled as operator valued distributions. Conceptually, this corresponds to a type of generalized function that maps each point in space time to a given operator. This mathematical abstraction is known to be leaky though, since multiplying two distributions is often an ill-specified operation, making it difficult to model interacting quantum fields.
\end{document}